\begingroup
\thispagestyle{empty}
\begin{tikzpicture}[remember picture,overlay]
\coordinate [below=8 cm] (midpoint) at (current page.north);
\node at (current page.north west)
{\begin{tikzpicture}[remember picture,overlay]
\node[anchor=north west,inner sep=0pt] at (0,0) {\includegraphics[width=\paperwidth]{epicstyle/background}}; % Background image
\draw[anchor=north] (midpoint) node [fill=ocre!30!white,fill opacity=0.6,text opacity=1,inner sep=1cm]{\Huge\centering\bfseries\sffamily\parbox[c][][t]{\paperwidth}{\centering GOSBook\\[15pt] % Book title
%{\Large }\\[20pt] % Subtitle
{\huge Didenko Andre}}}; % Author name
\end{tikzpicture}};
\end{tikzpicture}
\vfill
\endgroup

\newpage
\thispagestyle{empty}
\begingroup
\newgeometry{margin=1.6cm, top=1.2cm, footskip = 1 cm, headheight=11pt,headsep=0.2cm}
\epigraph{Умереть не страшно, страшно не жить.}{Виктор Гюго}\centering

\settowidth{\unitlength}{\LARGE\scshape Московский Физико-Технический Институт}
\vspace{-2.2\baselineskip}
\begin{figure}[!h]
\center
\includegraphics[width=0.27\textwidth]{pictures/MIPT2}
\end{figure}
\vspace*{0.5\baselineskip}
{\LARGE\scshape Московский Физико-Технический Институт}\\[0.7\baselineskip]
\rule{\unitlength}{1.6pt}\vspace*{-\baselineskip}\vspace*{2pt}
\rule{\unitlength}{0.4pt}\\[1.4\baselineskip]
{\Huge Подготовка к ГОСу по Матану}\\[\baselineskip]
{\large \itshape Sapere aude}\\[0.8\baselineskip]
\rule{\unitlength}{0.4pt}\vspace*{-1.5\baselineskip}\vspace{3.2pt}
\rule{\unitlength}{1.6pt}\\[\baselineskip]
{\Large\scshape Диденко А.А. \\\vspace*{0.2\baselineskip}  \normalsize студент 311 группы}\par
\vspace*{\baselineskip}  

%%%%%%%%%%%%%%%%%%%%%%%%%%%%%%%%%%%%%%%%%%%%%%%%%%%%%%%%%%%%%%%%%%%%%%%%%%%%%%%%%%%%%%%%%%%%%%%%%%%%%%%%%%%%%%%%%%%%%%%%%%%%%%%%%%%%%%%%%4
{\LARGE\scshape Редакция от \twodigit\day.\twodigit\month.\the\year \;(\currenttime)}\par % при печати убрать.

\bigskip
\vfill

\begin{flushleft}
$\bullet$ По любым вопросам касательно книги, а также если Вы хотите помочь в написании материала, просьба обращаться в личные сообщения: 

\vspace*{0.2\baselineskip}

\qquad\textcolor{Purplemountainmajesty}{$\blacktriangleright$}\href{https://vk.com/didenko_andre}{\large$vk.com/didenko\_andre$}

\vfill

$\bullet$ Актуальная версия доступна по ссылкам:

\vspace*{0.2\baselineskip}

\qquad\textcolor{Purplemountainmajesty}{$\blacktriangleright$}\href{http://latex.aslushnikov.com/compile?git=https://github.com/DidenkoAndre/GOS_book&target=_main.tex}{\large \textcolor{Purplemountainmajesty}{$latex.aslushnikov.com/...$}}

\vspace*{0.2\baselineskip}

\qquad\textcolor{Purplemountainmajesty}{$\blacktriangleright$}\href{https://www.dropbox.com/sh/7e5mfj8q68o2ipp/AAD8XvpZhiJzFbEh_IeH305ia?dl=0&preview=GOSBook.pdf}{\large \textcolor{Purplemountainmajesty}{$dropbox.com/...$}}

\vfill

$\bullet$ \textbf{Внимание.} Для Вас работает Telegram-бот \href{https://t.me/GOSBook_Bot}{\textcolor{PersianBlue}{$@GOSBook\_Bot$}}. Оформив подписку\footnotemark{}, Вы будете получать свежую версию автоматически!

\footnotetext{Для оформления подписки используйте команду \textsf{/subscribe}.}

\vfill

$\bullet$ Дополнительные материалы, полезные при подготовке как к \textit{устному} экзамену, так и к \textit{письменному} доступны по ссылке:

\vspace*{0.2\baselineskip}

\qquad\textcolor{Purplemountainmajesty}{$\blacktriangleright$}\href{https://drive.google.com/drive/u/0/folders/0BzuzEyNkpwYDYjVNcE0wa3hqWjA}{\large $drive.google.com/...$}

\vfill

$\bullet$ Вы всегда можете поддержать проект материально! 
Перевести средства на развитие проекта можно на карту {\large 5244 6873 8510 3816} \textsf{(MasterCard)}.

\bigskip
\vfill

\bigskip
\vfill

\textcolor{red}{
$\bullet$ Я крайне не рекомендую использовать данное пособие как \textit{\textbf{основную}} литературу во время подготовки к Государственному экзамену.
}
\end{flushleft}

\vfill
{\LARGE\scshape \LaTeX}\\[0.5\baselineskip]
{\LARGE\scshape 2016}\par
\restoregeometry
\endgroup
\newpage